% ==========================================================================
% GENERATED by scripts/split_paper.py. Do not edit this file.
%
% Source: imported_papers/grokking/iclr2027_conference_v2.tex
% Part 11: Controls and alternative explanations
%
% Edits belong upstream, in the published paper. To change what the
% chapter includes, change imports/grokking.toml and re-run
% the splitter.
% ==========================================================================
\section{Controls and alternative explanations}
Each part of this section closes one alternative reading of the results.

\subsection{Matched cosine controls on credit direction}
\label{app:mech}
Two simpler explanations of DFC's distributed circuit come before ours, and both are wrong. The first is that DFC is backpropagation with a hidden penalty attached, reaching the gradient by another route. It is not: throughout training its update points about $66$ degrees off the true gradient, a cosine of $0.4$ rather than $1$. Cosine is how we report misalignment below, with $1$ along the true gradient and $0$ at right angles.

The second is that misalignment of that size is itself what distributes a circuit. Two matched-cosine controls settle it (Table~\ref{tab:direction}). Both force backpropagation's update to a fixed cosine with the true gradient, and they differ only in whether the transverse part is redrawn each step. Redrawn, transverse components in random directions cancel and grokking survives to cosine $0.2$. Held fixed, grokking is gone by cosine $0.7$, and by $0.4$ the network memorizes at feedback alignment's own per-neuron periodicity. The size of the misalignment is not what decides.

Consistency comes closer and is still not the answer, and DFC is what rules it out. We measure persistence by removing the update's projection onto the true gradient and correlating the remaining direction with itself $4000$ epochs later, so $1$ means one direction held. DFC's component is as persistent as the fixed bias control, $0.97$ against $0.99$, and it groks where the control memorizes. Nor does any scalar alignment number order the rules. Measured at the input embedding alone, K-FAC sits at cosine $0.98$ and memorizes, while DFC sits at $0.72$, below the $0.78$ at which feedback alignment memorizes, and groks.

What separates DFC is that its push is tied to its own error: the controller signal shrinks as the output reaches the target, so at the solution there is nothing left to push with, where a fixed bias never shrinks. The error-gated control isolates exactly that. Both arms apply one fixed bias per seed at the same initial cosine, and the gated arm additionally scales the transverse part by the current loss, fading it to $0.011$ of its initial size while its direction never changes. The constant arm groks $0/3$ at every cosine below $1$, and the gated arm $3/3$, at test $0.984$ to $0.995$. Fading a push that would otherwise destroy grokking restores it completely, which is what Proposition~\ref{prop:persist} requires.

\begin{table}[!ht]
\caption{Persistent misalignment, not misalignment itself, is what destroys grokking. Three matched-cosine controls on credit direction (modular addition, $p=31$, three seeds), all forcing the update to a fixed cosine with the true gradient at matched magnitude. Fresh noise is redrawn each step, the persistent bias holds one direction, and the error-gated arm holds the same direction scaled by the current loss, so it vanishes at the solution. DFC operates near cosine $0.4$ with an error-gated residual, and groks.}
\label{tab:direction}
\begin{center}
\small
\begin{tabular}{lccccc}
\toprule
Cosine to gradient & $1.0$ & $0.7$ & $0.5$ & $0.4$ & $0.2$ \\
\midrule
i.i.d.\ noise: grok rate       & $3/3$ & $3/3$ & $3/3$ & $3/3$ & $3/3$ \\
i.i.d.\ noise: spectral entropy & $0.85$ & $0.83$ & $0.85$ & $0.85$ & $0.84$ \\
i.i.d.\ noise: per-neuron      & $0.64$ & $0.57$ & $0.56$ & $0.56$ & $0.54$ \\
\midrule
persistent bias: grok rate     & $3/3$ & $0/3$ & $0/3$ & $0/3$ & $0/3$ \\
persistent bias: test          & $0.998$ & $0.607$ & --- & $0.000$ & $0.000$ \\
persistent bias: per-neuron    & $0.64$ & $0.44$ & $0.26$ & $0.23$ & $0.22$ \\
\midrule
error-gated bias: grok rate    & $3/3$ & $3/3$ & --- & $3/3$ & $3/3$ \\
error-gated bias: test         & $0.998$ & $0.984$ & --- & $0.995$ & $0.991$ \\
error-gated bias: per-neuron   & $0.63$ & $0.78$ & --- & $0.77$ & $0.75$ \\
\bottomrule
\end{tabular}
\end{center}
\end{table}

\paragraph{The criterion measured on trained rules.} The controls impose the criterion, and this measurement reads it off the rules as trained. At checkpoints through training we compare each rule's applied update $\vu$ with the true gradient $\vg$ at the same parameters, recording the transverse part $\vu_\perp$ in units of $\|\vg\|$ and the autocorrelation of its direction (Table~\ref{tab:survival}). Backpropagation and the knob are excluded by construction, since a positive reweighting's transverse part vanishes wherever the gradient does and Proposition~\ref{prop:persist}(i) cannot displace its critical points. The exclusion is itself a prediction, that every positive reweighting groks, which the knob, LARS and Nero all do, and the sign rules are scored by their convergence theory instead (Appendix~\ref{app:theory}).

\begin{table}[!ht]
\caption{Rules that memorize keep their transverse component at the solution, while DFC collapses its own and groks. The transverse component of each rule's update, measured against the true gradient at the same parameters (modular addition, $p=31$, three seeds). Early and late are the first and last quarter of checkpoints, late/early the ratio of the raw $\|\vu_\perp\|$, and \emph{survives} means a component of held direction is still there at the solution, the memorizing outcome. Backpropagation and the knob are excluded as positive reweightings.}
\label{tab:survival}
\begin{center}
\small
\begin{tabular}{lcccccc}
\toprule
Rule & Test & $\|\vu_\perp\|/\|\vg\|$ early & late & Late/early & Direction lag & Survives \\
\midrule
Backpropagation & $0.996$ & $0.000$ & $0.000$ & --- & $0.03$ & excluded \\
Knob ($\beta{=}1$) & $0.976$ & $0.132$ & $0.090$ & $0.57$ & $0.57$ & excluded \\
\midrule
Fixed bias & $0.000$ & $0.917$ & $0.917$ & $0.87$ & $0.99$ & yes \\
DFA & $0.000$ & $0.344$ & $0.234$ & $0.61$ & $0.94$ & yes \\
FA (layerwise) & $0.000$ & $0.340$ & $0.235$ & $0.60$ & $0.92$ & yes \\
\midrule
Fresh noise & $0.986$ & $0.917$ & $0.917$ & $0.77$ & $-0.00$ & no \\
Error gated bias & $0.995$ & $0.092$ & $0.026$ & $0.07$ & $0.99$ & no \\
DFC & $1.000$ & $0.410$ & $0.038$ & $0.16$ & $0.74$ & no \\
\bottomrule
\end{tabular}
\end{center}
\end{table}

Among the rules it covers the separation is complete: every rule that memorizes keeps a transverse component of substantial size whose direction persists at $0.92$ or above, and every rule that groks fails that description. A component destroys the circuit only when it keeps both its size at the solution and one direction. DFC keeps the direction, persistent at $0.74$, and loses the size, falling by a factor of eleven as the controller signal vanishes. Fresh noise keeps the size and holds no direction. Both grok, which places DFC with the tolerated cases rather than with the fixed bias it otherwise resembles.

\subsection{The same credit under two optimizers}
\label{app:consumer}
Everything above compares rules that hand their credit to one shared optimizer. Damped Gauss-Newton and K-FAC instead apply their own step, and both precondition the gradient by curvature. Both are positive definite reweightings, so their update vanishes wherever the gradient does, and Proposition~\ref{prop:persist}(i) therefore cannot displace their solutions however far from the gradient they point.

Applied with a raw step both memorize, and consumed by the same optimizer every other rule in the paper uses both grok (Table~\ref{tab:consumer}). The credit is identical in each pair and only the optimizer that applies it differs, so the failure was never a property of the credit, which is why K-FAC appears in the bottom block of Table~\ref{tab:twoaxis} with the configuration it was tested at rather than as a rule that cannot grok. Where they grok, their circuits also land by grain: Gauss-Newton stays concentrated at participation $0.540$ to $0.619$ across dampings, and K-FAC, the one rule of the pair that whitens per unit, lands at $0.790$, outside the concentrated range of every coarse grained rule we ran. The composition also runs the other way: handing the knob's equalized credit to Lion instead of the shared optimizer still lands distributed, at circuit participation $0.984\,(0.003)$ and test $0.974\,(0.010)$ over three seeds, where Lion on its own credit sits concentrated near $0.5$, so the per-unit rescaling travels with the credit whoever consumes it.

The same dependence shows up in the fragility we report elsewhere: DFC's narrow window (Appendix~\ref{app:dfcwindow}) and sign-SGD's sensitivity to its learning rate are the same phenomenon, a transverse component that is small enough at one operating point and not at another.

\begin{table}[!ht]
\caption{The same credit under two optimizers gives opposite outcomes, for both curvature rules. Test accuracy on modular addition ($p=31$, three seeds). A raw step is the rule applying its own update directly, and the shared optimizer is the one every other rule in the paper uses. Damping is the term added to the curvature matrix before inversion.}
\label{tab:consumer}
\begin{center}
\small
\begin{tabular}{llcc}
\toprule
Rule & Damping & Raw step & Shared optimizer \\
\midrule
Gauss-Newton & $10^{-2}$ & $0.000$ & $0.990$ \\
             & $10^{-1}$ & $0.000$ & $0.994$ \\
             & $1$       & ---     & $0.998$ \\
\midrule
K-FAC        & $10^{-1}$ & $0.001$ & $1.000$ \\
\bottomrule
\end{tabular}
\end{center}
\end{table}

\subsection{Optimizer and regularizer baselines}
\label{app:baselines}
Three familiar mechanisms could in principle produce a distributed circuit: an adaptive optimizer's preconditioning, ordinary regularization, and hyperparameter tuning. None of them produces ours (Table~\ref{tab:baselines}).

Adam equalizes per parameter, not per unit, so it neither produces this effect nor cancels it. Rerunning the sweep under plain SGD with momentum reproduces both endpoints, concentrated at $\beta=0$ (participation $0.53$, deletion tolerance $0.60$) and distributed at $\beta=1$ ($0.99$, $0.76$), and Adam's own effective per-unit credit, the participation of $m/\sqrt{v}$, reaches only $0.66$ at $\beta=0$ against the raw gradient's $0.59$ and $0.99$ at $\beta=1$, so it passes the knob's equalization through rather than supplying or absorbing it. Dropout does reach a distributed circuit, participation up to $0.99$, which is expected since it is a redundancy pressure, but it gets there by changing the objective, its solution drifting to higher margin ($5.3 \to 6.8$) and weight norm ($26.8 \to 30.2$) as the rate grows. Several routes lead to a distributed circuit, and the credit rescaling is the one that reaches it at backpropagation's own loss, margin and weight norm, changing nothing about the problem being solved.

Backpropagation cannot reach the distributed end at any setting of its own hyperparameters either. Sweeping three learning rates crossed with three weight decays, six of the nine grok and all six land concentrated, spanning only $0.530\,(0.008)$ to $0.626\,(0.012)$, short of even the midpoint between the two ends (Table~\ref{tab:bpgrid}). The three failures are at the highest weight decay, where test accuracy falls to $0.21$ to $0.36$, and a memorizing network's participation is not a grokked circuit. The grid also reproduces the agreement between credit and circuit six more times, with credit from $0.523$ to $0.613$ predicting circuit participation to a mean absolute deviation of $0.009$. The knob settings between the two ends fill the curve monotonically, $\beta=0.5$ at credit $0.99$ and circuit $0.98\,(0.00)$ with test $0.985$, and $\beta=0.75$ at $1.00$ and $0.99\,(0.01)$ with test $0.977$.

\begin{table}[!ht]
\caption{The credit-magnitude endpoint is not an Adam or dropout artifact ($p=31$,
three seeds). SGD+momentum reproduces the $\beta$ endpoints, Adam's effective per-unit credit stays concentrated, and dropout reaches a distributed circuit only by raising the margin and norm, which the knob does not.}
\label{tab:baselines}
\begin{center}
\small
\begin{tabular}{lccccc}
\toprule
Setting & Circuit PR & Random del & Margin & Norm & Note \\
\midrule
SGD+mom, $\beta{=}0$ & $0.53$ & $0.60$ & --- & --- & concentrated \\
SGD+mom, $\beta{=}1$ & $0.99$ & $0.76$ & --- & --- & distributed (as Adam) \\
Adam raw grad ($\beta{=}0$) & $0.59$ & --- & --- & --- & credit participation \\
Adam effective ($\beta{=}0$) & $0.66$ & --- & --- & --- & $m/\sqrt{v}$ credit participation \\
Adam effective ($\beta{=}1$) & $0.99$ & --- & --- & --- & Adam propagates the knob \\
LARS (per layer) & $0.63$ & $0.67$ & --- & --- & groks, stays concentrated \\
Dropout $0.1$ & $0.95$ & $0.79$ & $5.3$ & $26.8$ & objective changed \\
Dropout $0.5$ & $0.99$ & $0.82$ & $6.8$ & $30.2$ & objective changed \\
Credit knob $\beta{=}1$ & $0.99$ & $0.77$ & $5.2$ & $26.6$ & objective preserved \\
\bottomrule
\end{tabular}
\end{center}
\end{table}

\begin{table}[!ht]
\caption{Every setting that groks lands concentrated, and the applied credit tracks the grokked circuit throughout. Backpropagation under its own tuning ($p=31$, three seeds, standard deviations in parentheses). The three highest weight decay settings memorize, so their participation does not describe a grokked circuit.}
\label{tab:bpgrid}
\begin{center}
\small
\begin{tabular}{lccccc}
\toprule
Learning rate & Weight decay & Test & Credit PR & Circuit PR & Groks \\
\midrule
$3\times10^{-4}$ & $10^{-3}$        & $0.998$ & $0.613$ & $0.626\,(0.012)$ & yes \\
$3\times10^{-4}$ & $3\times10^{-3}$ & $0.983$ & $0.560$ & $0.571\,(0.027)$ & yes \\
$10^{-3}$        & $10^{-3}$        & $0.997$ & $0.604$ & $0.613\,(0.016)$ & yes \\
$10^{-3}$        & $3\times10^{-3}$ & $0.993$ & $0.523$ & $0.530\,(0.008)$ & yes \\
$3\times10^{-3}$ & $10^{-3}$        & $0.993$ & $0.589$ & $0.594\,(0.024)$ & yes \\
$3\times10^{-3}$ & $3\times10^{-3}$ & $0.988$ & $0.546$ & $0.553\,(0.022)$ & yes \\
\midrule
$3\times10^{-4}$ & $10^{-2}$        & $0.357$ & $0.474$ & $0.544\,(0.058)$ & no \\
$10^{-3}$        & $10^{-2}$        & $0.229$ & $0.367$ & $0.434\,(0.157)$ & no \\
$3\times10^{-3}$ & $10^{-2}$        & $0.212$ & $0.417$ & $0.505\,(0.338)$ & no \\
\bottomrule
\end{tabular}
\end{center}
\end{table}

\subsection{DFC's grokking window}
\label{app:dfcwindow}
DFC's grokking of modular addition is real but sensitive to hyperparameters, as grokking generally is. Table~\ref{tab:dfcwindow} sweeps weight decay and the controller step size at $p=31$. DFC groks on both seeds at weight decay $10^{-3}$ with controller step size in $\{0.2, 0.3\}$, converging to a distributed circuit, and outside this window it underfits or collapses. Our DFC results use weight decay $10^{-3}$ and controller step size $0.3$, inside the window. The window is narrower than backpropagation's, so DFC's role in this paper is to establish that a controller based rule reaches the distributed end at all.

\begin{table}[!ht]
\caption{DFC groks only at weight decay $10^{-3}$ with a large enough controller step. Final test accuracy on modular addition ($p=31$, two seeds) across weight decay and controller step size, where bold cells grok on both seeds.}
\label{tab:dfcwindow}
\begin{center}
\begin{tabular}{lccc}
\toprule
 & step size $0.1$ & step size $0.2$ & step size $0.3$ \\
\midrule
weight decay $10^{-3}$        & $0.68$ & $\mathbf{1.00}$ & $\mathbf{1.00}$ \\
weight decay $3\times10^{-3}$ & $0.01$ & $0.01$ & $0.74$ \\
weight decay $10^{-2}$        & $0.01$ & $0.01$ & $0.01$ \\
\bottomrule
\end{tabular}
\end{center}
\end{table}
